\subsection{KRULL 整环的本性赋值}

设 A 是 Krull 整环，K 是其分式域。第 3 节的公式 (4) 定义的赋值（对 \(x \in \mathbf{K}^*\)）称为 K（或 A）的本性赋值。

我们在定理 2 的证明过程中已经指出，赋值 \(v_{\mathbb{P}}\) 满足定义 3 的性质 \((\mathrm{AK}_{\mathbb{I}}), (\mathrm{AK}_{\mathbb{II}})\) 和 \((\mathrm{AK}_{\mathbb{III}})\)。此外，这些离散赋值 \(v_{\mathbb{P}}\) 是赋范的：对于每个极端除子 \(\mathbf{P} \in \mathbf{P}(\mathbf{A})\)，有 \(\mathbf{P} < 2\mathbf{P}\)，因而若 \(\mathfrak{a}\) 和 \(\mathfrak{b}\) 是分别对应于 \(\mathbf{P}\) 和 \(2\mathbf{P}\) 的除性理想，则 \(\mathfrak{a} \supset \mathfrak{b}\) 且 \(\mathfrak{a} \neq \mathfrak{b}\)；对于 \(x \in \mathfrak{a} - \mathfrak{b}\)，有 \(\operatorname{div}(x) \geqslant \mathbf{P}\) 且 \(\operatorname{div}(x) \not\geqslant 2\mathbf{P}\)，从而 \(v_{\mathbb{P}}(x) = 1\)，这就证明了断言。

命题 5. 设 A 是 Krull 整环，K 是其分式域，\((v_{\mathbf{P}})_{\mathbf{P} \in \mathbf{P}(\mathbf{A})}\) 是其本性赋值族。设 \((n_{\mathbf{P}})_{\mathbf{P} \in \mathbf{P}(\mathbf{A})}\) 是一个有理整数族，除有限个指标外其余值均为零。则集合 \(x \in \mathbf{K}\)，满足 \(v_{\mathbf{P}}(x) \geqslant n_{\mathbf{P}}\) 对所有 \(\mathbf{P} \in \mathbf{P}(\mathbf{A})\) 都成立，正是除性理想 \(\mathfrak{a}\)（\(\mathbf{A}\) 的），且 \(\operatorname{div} \mathfrak{a} = \sum_{\mathbf{P} \in \mathbf{P}(\mathbf{A})} n_{\mathbf{P}} \cdot \mathbf{P}\)。

令 \(x \in K^*\)。\(x \in a\) 的充要条件是 \(Ax \subset a\)，从而等价于 \(\operatorname{div}(x) \geqslant \operatorname{div} a\)，进而由 (4) 等价于 \(v_{\mathbf{P}}(x) \geqslant n_{\mathbf{P}}\) 对所有 \(\mathbf{P} \in \mathbf{P}(\mathbf{A})\) 都成立。

命题 6. 设 A 是 Krull 整环，K 是其分式域，\((v_{\iota})_{\iota \in I}\) 是 K 上满足定义 3 中性质的赋值族，\(\mathbf{A}_{\iota}\) 是 \(v_{\iota}\) 的赋值环。设 S 是 A 的不含 0 的乘性子集，J 是指标 \(\iota \in I\) 中满足 \(v_{\iota}\) 在 S 上为零者的集合。则 \(\mathrm{S}^{-1}\mathrm{A} = \bigcap_{\iota \in I}\mathrm{A}_{\iota}\)；特别地，\(\mathrm{S}^{-1}\mathrm{A}\) 是 Krull 整环。

记 \(\mathbf{B} = \bigcap_{\iota \in \mathbf{J}} \mathbf{A}_{\iota}\)。则 \(\mathbf{S}^{-1} \subset \mathbf{B}\) 且 \(\mathbf{A} \subset \mathbf{B}\)，从而 \(\mathbf{S}^{-1}\mathbf{A} \subset \mathbf{B}\)。反之，设 \(x \in \mathbf{B}\)。记 \(\mathbf{J}'\) 为指标 \(\iota\) 中满足 \(v_{\iota}(x) < 0\) 的有限集合。若 \(\iota \in \mathbf{J}'\)，则 \(x \notin \mathbf{A}_{\iota}\)，从而 \(\iota \notin \mathbf{J}\)，因而存在 \(s_{\iota} \in \mathbf{S}\) 使 \(v_{\iota}(s_{\iota}) > 0\)。取 \(n(\iota)\) 为一个整数 \(>0\)，使得 \(v_{\iota}(s_{\iota}^{n(\iota)}x) \geqslant 0\)；记 \(s = \prod_{\iota \in \mathbf{J}} s_{\iota}^{n(\iota)}\)。于是 \(v_{\iota}(sx) \geqslant 0\) 对所有 \(\iota \in \mathbf{I}\) 都成立，从而 \(sx \in \mathbf{A}\) 且 \(x \in \mathbf{S}^{-1}\mathbf{A}\)。因此 \(\mathbf{B} = \mathbf{S}^{-1}\mathbf{A}\)。

推论 1. 设 P 是 A 的一个极端除子，p 是相应的除性理想。则 p 是素理想，\(v_{\mathbf{P}}\) 的赋值环是 \(A_{\mathfrak{p}}\)，且 \(v_{\mathbf{P}}\) 的剩余域可与 A/p 的分式域相识别。

令 \(S = A - p\)。由命题 5，\(v_{P}\) 在 \(S\) 上为零，且 \(>0\) 在 \(p\) 上成立。因此，\(p\) 是 \(A\) 与 \(v_{P}\) 的理想的交，因而是素理想。另一方面，对于每个极端除子 \(Q \neq P\)，都有 \(Q \not\geq P\)，所以除性理想 \(q\)（对应于 \(Q\)）不包含于 \(p\)；于是 \(q \cap S \neq \varnothing\)，从而由命题 5，\(v_{Q}\) 在 \(S\) 上不为零。由此，推论由命题 6 及第 II 章，§ 3，第 1 节，命题 3 得出。

推论 2. 设 A 是 Krull 整环，K 是其分式域，\((v_{\iota})_{\iota\in I}\) 是满足定义 3 中性质的赋值族。则 A 的每个本性赋值都等价于某个 \(v_{\iota}\)。

设 P 是 A 的一个极端除子，p 是相应的除性理想。由推论 1、命题 5、引理 1 以及第 3 节定理 2 的证明中的断言 (1)，存在 \(\iota \in I\)，使得赋值环 \(A_{\iota}\)（属于 \(v_{\iota}\)）包含赋值环 \(A_{p}\)（属于 \(v_{P}\)）。由于 \(v_{\iota}\) 和 \(v_{P}\) 的高度都是 1，它们因而等价（第 VI 章，§4，第 5 节，命题 6）。

命题 7. 设 A 是 Krull 整环，\((v_{\mathbb{P}})_{\mathbb{P} \in \mathbb{P}(\mathbb{A})}\) 是其本性赋值族，且 \(\mathfrak{a} \in \mathbf{I}(\mathbf{A})\)。则 div \(\mathfrak{a}\) 中 P 的系数是 \(\inf_{y \in \mathfrak{a}} (v_{\mathbb{P}}(y))\)。若 \(\mathfrak{p}\) 是对应于极端除子 P 的除性素理想，则 \(\mathfrak{a}\mathrm{A}_{\mathfrak{p}} = \tilde{\mathfrak{a}}\mathrm{A}_{\mathfrak{p}}\)。

由于 \(\mathfrak{a} = \sum_{y \in \mathfrak{a}} Ay\)，第 1 节的命题 2 (b) 表明 \(\text{div}(\mathfrak{a}) = \inf_{y \in \mathfrak{a}} (\text{div}(Ay))\)，从而得到第一个断言。第二个断言立即成立，因为 \(\text{div} \tilde{\mathfrak{a}} = \text{div} \mathfrak{a}\)，且 \(A_{\mathfrak{p}}\) 是离散赋值 \(v_{\mathfrak{p}}\) 的赋值环。

命题 8. 设 A 是整闭的诺特整环。

\begin{enumerate}
\def\labelenumi{(\alph{enumi})}
\tightlist
\item
  设 \(\mathbf{P}\) 是 \(\mathbf{A}\) 的一个极值除子，\(\mathfrak{p}\) 是相应的除性素理想；
\end{enumerate}

对于 \(n \in \mathbf{N}\)，令 \(\mathfrak{p}^{(n)} = \mathfrak{p}^n\mathbf{A}_{\mathfrak{p}} \cap \mathbf{A}\)；则 \(\mathfrak{p}^{(n)}\) 是 \(x \in \mathbf{A}\) 的集合，其中 \(v_{\mathfrak{p}}(x) \geqslant n\)，并且是一个 \(\mathfrak{p}\)-准素理想。

\begin{enumerate}
\def\labelenumi{(\alph{enumi})}
\setcounter{enumi}{1}
\tightlist
\item
  设 \(\mathfrak{a}\) 是除性整理想，\(n_1\mathrm{P}_1 + \dots + n_r\mathrm{P}_r\) 是 \(\mathfrak{a}\) 的除子（其中 \(\mathbf{P}_i\) 是互异的极端除子），\(\mathfrak{p}_i\) 是对应于 \(\mathbf{P}_i\) 的除性素理想。
\end{enumerate}

则 \(\mathfrak{a} = \bigcap_{i=1}^{r} \mathfrak{p}_i^{(n_i)}\) 是 \(\mathfrak{a}\) 的唯一既约准素分解，且 \(\mathfrak{p}_i\) 不是浸入的。

由命题 6 的推论 1，关系 \(x \in \mathfrak{p}^n\mathrm{A}_{\mathfrak{p}} = (\mathfrak{p}\mathrm{A}_{\mathfrak{p}})^n\) 等价于 \(v_{\mathfrak{p}}(x) \geqslant n\)；另一方面，由于 \(\mathrm{A}_{\mathfrak{p}}\) 是离散赋值环，\((\mathfrak{p}\mathrm{A}_{\mathfrak{p}})^n\) 是 \((\mathfrak{p}\mathrm{A}_{\mathfrak{p}})\)-准素理想（第 IV 章，§ 2，第 1 节，例 4），从而 \(\mathfrak{p}^{(n)}\) 是 \(\mathfrak{p}\)-准素理想（第 IV 章，§ 2，第 1 节，命题 3）；这就证明了 (a)。命题 5 显然证明了 \(\mathfrak{a} = \bigcap_{i=1}^{r} \mathfrak{p}_i^{(n_i)}\)。由于 \(\mathfrak{p}_i \not\subset \mathfrak{p}_j\) 对 \(i \neq j\) 成立，这个准素分解是既约的：事实上，若 \(\mathfrak{p}_i^{(n_i)} \supset \bigcap_{j \neq i} \mathfrak{p}_j^{(n_j)} \supset \prod_{j \neq i} \mathfrak{p}_j^{(n_j)}, \mathfrak{p}_i\) 将包含某个 \(\mathfrak{p}_j\)，其中 \(j \neq i\)（第 II 章，§ 1，第 1 节，命题 1）。唯一性由第 IV 章，§ 2，第 3 节，命题 5 得出。

